% !TeX root = main.tex
\newpage
\section{Boundary-Space Lighting}

\subsection{Purpose}

Boundary-space lighting is the active surface-lighting direction for
\texttt{LightingBall}.  Its purpose is to let photons light part of an
analytic or implicit surface without turning the moving photon point into the
visible lighting result.

The current model separates responsibilities:
\begin{itemize}
    \item particles carry motion, identity, and collision/debug state;
    \item materials provide source color and debug visibility settings;
    \item boundary markers make analytic surface regions discoverable;
    \item \texttt{BoundaryLightRecord} stores persistent surface color;
    \item the renderer displays that persistent color on the surface.
\end{itemize}

This keeps photons as transport/debug particles while allowing a surface region
to remain lit after contact.

\subsection{Cell-Owned Surface State}

The first Vulkan proxy-kernel implementation stored one lighting record per
boundary particle.  That proved the contact-to-light path, but sparse boundary
markers produced point-like circles, fan-out, and visible gaps.  The active
Vulkan model now stores lighting by spatial cell instead.

Conceptually:
\begin{verbatim}
cell contains a small surface region
photon hits anywhere in that region
BoundaryLight[cell_id] stores the material color
renderer colors the analytic surface from the owning cell
\end{verbatim}

The boundary marker keeps its base material identity and broad-phase locality
role.  The visible light state is separate resource state:
\begin{verbatim}
material_id                         = base surface material
BoundaryLight[cell_id].rgb_valid.rgb = persistent lit color
BoundaryLight[cell_id].rgb_valid.w   = valid flag
\end{verbatim}

\subsection{Photon Hit Rule}

The first rule is intentionally simple:
\begin{verbatim}
photon hits surface
cell_id = address of contacted surface point
BoundaryLight[cell_id].rgb_valid = contacted material color, valid=true
\end{verbatim}

The color does not clear every compute frame.  It persists until another photon
updates the same cell.  This is a color-state rule, not an energy-accumulation
rule.

Because the first model sets color instead of adding energy, no atomic add is
required.  If two photons write the same cell in the same compute dispatch, the
first implementation accepts last-writer-wins behavior.  If later work needs
additive light transport, atomic accumulation can be introduced as a separate,
deliberate feature.

\subsection{Analytic Sphere Surface}

For \texttt{LightingBall}, photon contact is evaluated against the analytic
sphere, not against the marker particles.  When a photon contacts the sphere,
the compute shader writes the cell containing the analytic hit point:
\begin{verbatim}
hitPoint = LIGHTING_BALL_CENTER + contactNormal * LIGHTING_BALL_RADIUS
cell_id = CellAddress(hitPoint)
BoundaryLight[cell_id].rgb_valid = materialColor, valid=true
\end{verbatim}

The lighting sphere render resource uses the existing sphere mesh only as a
canvas.  During \texttt{ResourceLightingSphere::Create}, each unit-sphere
vertex is mapped to the configured world-space sphere and assigned an owning
cell address:
\begin{verbatim}
sphereNormal = normalize(unit sphere vertex)
worldPos = LIGHTING_BALL_CENTER + sphereNormal * LIGHTING_BALL_RADIUS
surfaceCellId = CellAddress(worldPos)
\end{verbatim}

The vertex shader passes \texttt{surfaceCellId} to the fragment shader as a
flat value.  The fragment shader reads:
\begin{verbatim}
BoundaryLightRecord record = BoundaryLight[surfaceCellId];
\end{verbatim}

If the record is valid and belongs to the lighting-ball surface, the fragment
uses the stored color at full alpha.  If it is invalid, the fragment is
discarded.  This gives solid cell-owned surface regions without angular
fan-out, point-sprite circles, or proxy-color averaging.

\subsection{Point-List And Surface Rendering}

The ordinary particle/debug pass still uses:
\begin{verbatim}
inputAssembly.topology = VK_PRIMITIVE_TOPOLOGY_POINT_LIST;
\end{verbatim}

That pass remains useful for photons and particle diagnostics.  It is not the
final surface canvas.  \texttt{ParticleLighting} now also has a
lighting-specific surface pass for \texttt{LightingBall}.  The mesh used by
that pass is render geometry only; collision remains analytic, and the
generated boundary markers remain broad-phase locality records.

\subsection{Vulkan Implementation}

The Vulkan version uses the existing record name:
\begin{verbatim}
struct BoundaryLightRecord {
    vec4 rgb_valid;
    vec4 normal_material;
    uvec4 ids;
};
\end{verbatim}

The buffer is owned by \texttt{ResourceLighting.cpp}, sized as
\texttt{CellAryW * CellAryH * CellAryL}, and bound at descriptor binding 8
only for the \texttt{ParticleLighting} application.  Basic particle
applications remain on \texttt{FPM.*} and do not see this resource.

\texttt{FPML.comp} is the lighting compute schedule.  It deposits contacted
material color into \texttt{BoundaryLight[cell\_id]} with persistent
last-writer-wins behavior.  \texttt{BoundaryLight.vert} passes the baked
surface cell id to \texttt{BoundaryLight.frag}, and the fragment shader reads
the cell-owned color directly.  Since the fragment shader reads the lighting
SSBO, the \texttt{ParticleLighting} compute-finished wait includes the fragment
shader stage.

\subsection{Python Reference}

Python remains the behavioral reference for photon contact and material color
selection.  Its renderer may still use existing boundary markers as a coarse
visual approximation, but the Vulkan target is cell-owned analytic surface
color.  The Python model should not reintroduce the retired
\texttt{BoundaryLightSample}/\texttt{BoundaryLightState} current/filter path or
make the photon point itself the final visible surface lighting.

\subsection{Deferred Spectral Lighting}

Spectral lighting is not working in the current cell-owned Vulkan milestone.
The next lighting work should decide what additional state is required beyond
\texttt{rgb\_valid}.  Candidate future state includes photon wavelength or
band, material spectral response, accumulated direction/coherence, or a
separate debug view for spectral deposits.  That work should extend the
cell-owned \texttt{BoundaryLightRecord} path deliberately rather than reviving
the old proxy-kernel or per-frame filtered-light buffers.

\subsection{Decision}

The active Vulkan path is:
\[
    \textrm{photon hit}
    \rightarrow
    \textrm{persistent cell-owned surface color}
    \rightarrow
    \textrm{analytic surface fragment color}.
\]

This preserves the useful lesson from the earlier surface-cell idea: lighting
must belong to surface space.  The current implementation achieves that without
making the visible boundary particles carry the final surface color.
